%%% PREAMBLE %%%

% Strict mode

\RequirePackage[l2tabu, orthodox]{nag} % Warn when using deprecated constructs

% Document class and packages
\documentclass[10pt,a4paper,parskip,fleqn]{scrartcl}
\usepackage[a4paper,vmargin={30mm},hmargin={30mm}]{geometry} % Page margins
\usepackage[ngerman]{babel} % New German hyphenation (multilingual support)
\usepackage[utf8x]{inputenc} % Unicode support
\usepackage{graphicx} % Graphics support
\usepackage{enumitem}

% Font configuration
\usepackage[sc]{mathpazo} % Use palatino font
\usepackage[T1]{fontenc} % Use correct font encoding
\usepackage[babel=true]{microtype} % Micro-typographic optimizations
\addtokomafont{disposition}{\rmfamily} % Set palatino as heading font

% Section Titles
\usepackage{titlesec}
\titleformat{\section}{\Large\bfseries}{§ \thesection{} -- }{0em}{}

% Numbering
\renewcommand*{\labelenumi}{\thesection.\arabic{enumi}}
\renewcommand*{\labelenumii}{\labelenumi.\arabic{enumii}}

% Easier enumerations
\newcommand{\ol}{\begin{enumerate}[itemsep=-0.2em,topsep=-0.2em]}
\newcommand{\lo}{\end{enumerate}}
\newcommand{\li}{\item}


%%% TITLEPAGE %%%

\title{\Huge Vereinsstatuten}
\subject{open\textbackslash{}HSR}
\date{Stand: 23. Mai 2016}


%%% HEADER / FOOTER %%%

\usepackage{fancyhdr} % Fancy headers
\usepackage{lastpage} % Page numbering

% Capture title and author
\makeatletter
\let\Date\@date
\makeatother

% Fancy headers configuration
\pagestyle{fancy}
\fancyhead{} % Clear all header fields
\fancyhead[LO,L]{\bfseries Vereinsstatuten open\textbackslash{}HSR}
\fancyhead[RO,R]{\bfseries \Date{}}
\fancyfoot{} % Clear all footer fields
\fancyfoot[CO,C]{Seite \thepage}
\renewcommand{\headrulewidth}{0.3pt}
\renewcommand{\footrulewidth}{0pt}

% Better footer
\cfoot{Seite \thepage\ von \pageref{LastPage}}


%%% MAIN DOCUMENT %%%

\begin{document}

\begin{titlepage}
  \maketitle
  \thispagestyle{empty} % Don't start page numbers on this page
  \begin{center}
    \includegraphics[width=0.25\textwidth]{logo.png}\\

    \vfill

    \large Beschlossen an der GV vom 23. Mai 2016.

    \vspace{1.5cm}

    \begin{minipage}[t]{0.49\textwidth}
      \center
      \rule{5cm}{0.2mm}\\
      Fabian Hauser, Präsident
      \endcenter
    \end{minipage}
    \begin{minipage}[t]{0.49\textwidth}
      \center
      \rule{5cm}{0.2mm}\\
      Murièle Trentini, Kassierin
      \endcenter
    \end{minipage}
  \end{center}
\end{titlepage}


\section{Name und Sitz}

\ol
  \li Unter dem Namen \textit{open\textbackslash{}HSR} besteht ein Verein gemäss Art. 60ff
  ZGB.
  \li Der Verein ist religiös unabhängig und politisch neutral
  \li Der Sitz des Vereins ist Rapperswil-Jona. Für alle Rechtsstreitigkeiten
  gilt schweizerisches Recht.
\lo


\section{Vereinszweck}

\ol
  \li Der Verein hat folgende Ziele:
    \ol
      \li Förderung der Open Source-Kultur und von Open Source-Projekten an
      der HSR
      \li Interessensvertretung an der HSR für die Unterstützung
      alternativer Betriebssysteme
      \li Unterstützung von Einsteigern in der Nutzung von Open
      Source-Software im Hinblick auf das Studium an der HSR
    \lo
  \li Der Verein unterstützt die Ziele und Ideale des VSHSR. 
\lo


\section{Mitgliedschaft}

\ol
  \li Der Verein \textit{open\textbackslash{}HSR} steht allen aktiven oder ehemaligen
  Studenten, Assistenten, Dozenten oder sonstigen Angehörigen der Hochschule für
  Technik in Rapperswil-Jona (HSR) offen.
  \li Aufnahme von Mitgliedern
    \ol
      \li Natürliche Personen erlangen die Mitgliedschaft durch Mitteilung an
      ein Vorstandsmitglied und Bezahlung des Mitgliederbeitrages.
      \li Juristische Personen stellen einen schriftlichen Antrag an den
      Vorstand, welcher über die Aufnahme entscheidet.
      \li Stimm- und wahlberechtigt sind nur diejenigen Mitglieder, welche den
      Mitgliederbeitrag für das laufende Vereinsjahr bezahlt haben.
    \lo
  \li Austritt von Mitgliedern
    \ol
      \li Ein Austritt ist jederzeit möglich. Für das aktuelle Vereinsjahr
      bezahlte Mitgliederbeiträge werden nicht erstattet.
      \li Der Austritt hat nach einer der folgenden Möglichkeiten zu geschehen:
        \ol
          \li Schriftlich an den Vorstand (per Brief oder E-Mail).
          \li Mündlich an einer beschlussfähigen Vorstandssitzung oder einer
           Vereinsversammlung.
        \lo
      \li Wird nach Studienabschluss der Vereinsbeitrag während eines vollen
      Vereinsjahres nicht bezahlt, erlischt die Mitgliedschaft automatisch.
      \li Vom Verein ausgeschlossene oder ausgetretene Mitglieder haben keinen
      Anspruch auf Rückerstattungen jeglicher Art.
      \li Beim Austritt aus der HSR erlischt die Mitgliedschaft für Mitglieder
      ohne Vorstandsfunktion automatisch.
    \lo
  \li{Ausschluss von Mitgliedern}
    \ol
      \li Mitglieder des Vereins können ausgeschlossen werden, wenn sie dem
      Zweck des Vereins zuwiederhandeln, seinen Ruf oder denjenigen eines
      Vereinsmitglieds schädigen oder einen Vereinsbeschluss oder die Statuten
      missachten.
      \li Die Vereinsversammlung kann den Ausschluss mit einer
      Zweidrittelmehrheit (keine Berücksichtigung der Enthaltungen)
      beschliessen.
      \li Ein allfälliger Ausschluss muss innerhalb der üblichen Frist für
      Traktanden angekündigt werden.
      \li Vor der Beschlussfassung wird dem betreffenden Mitglied eine
      Einsprache gewährt.
    \lo
\lo


\section{Organisation}

\ol
  \li Die Organe des Vereins sind
    \ol
      \li die Vereinsversammlung
      \li der Vorstand
      \li die Arbeitsgruppen
    \lo
  \li Die Vereinsversammlung
    \ol
      \li findet mindestens einmal im Vereinsjahr statt.
      \li muss mindestens zwei Wochen im Voraus schriftlich angekündigt werden.
      \li wird vom Vorstand einberufen. Die Einberufung muss zudem zwingend
      erfolgen, wenn mindestens die Hälfte der Mitglieder dies verlangt.
      \li ist beschlussfähig, wenn entweder mindestens sieben Mitglieder oder
      ein absolutes Mehr der stimmberechtigten Mitglieder anwesend sind.
      \li genehmigt oder verwirft die Jahresrechnung.
      \li wählt den Präsidenten.
      \li wählt den restlichen Vorstand.
      \li entscheidet bei Stimmengleichheit mit dem Stichentscheid des
      Präsidenten.
      \li kann auch über nicht frühzeitig angekündigte Traktanden entscheiden,
      sofern dies von einer Mehrheit der anwesenden Mitglieder gutgeheissen
      wird.  Davon ausgeschlossen sind die Auflösung des Vereins, der Ausschluss
      von Mitgliedern, jegliche Zahlungen sowie Statutenänderungen.
    \lo
  \li Der Vorstand
    \ol
      \li besteht zwingend aus: Präsident, Vize-Präsident, Aktuar und Kassier.
      Zusätzlich können bis zu 4 weitere Vorstandsmitglieder bestimmt werden,
      die dieselben Rechte und Pflichten wie der restliche Vorstand haben.
      \li bestimmt seine Aufgaben und Ämterverteilung selbst, sofern nicht von
      der Mitgliederversammlung oder den Statuten anders festgelegt.
      \li ist beschlussfähig, sofern mindestens die Hälfte des Vorstands oder
      drei Vorstandsmitglieder anwesend sind.
      \li beruft die Vereinsversammlung ein.
      \li entscheidet über das Vereinsvermögen.
      \li entscheidet bei Stimmengleichheit mit dem Stichentscheid des
      Präsidenten. Bei dessen Abwesenheit liegt der Stichentscheid beim
      Vize-Präsidenten.
      \li für finanzielle Belange gilt für den Präsidenten und den
      Vizepräsidenten die Kollektiv\-unterschriftsberechtigung zu zweien, sowie
      für den Kassier Einzelunterschriftsberechtigung.
    \lo
  \li Die Arbeitsgruppen
    \ol
      \li werden bei Bedarf formlos gebildet und aufgelöst.
      \li erstatten dem Vorstand Bericht über ihre Tätigkeiten.
    \lo
\lo


\section{Finanzen}

\ol
  \li Der Verein arbeitet nicht auf einen finanziellen Gewinn hin.
  \li Der Mitgliederbeitrag wird jährlich erhoben.
  \li Die Höhe des Mitgliederbeitrags wird jährlich von der Vereinsversammlung
  festgelegt.
  \li Das Vereinsjahr beginnt mit dem Frühlings-Semester der HSR.
  \li Der Verein haftet allein mit dem Vereinsvermögen.
\lo


\section{Auflösung}

\ol
  \li Die Auflösung des Vereins kann durch die Vereinsversammlung mit einer
  Zweidrittelmehrheit beschlossen werden.
  \li Eine allfällige Auflösung muss innerhalb der üblichen Frist für Traktanden
  angekündigt werden.
  \li Sollte bei der Auflösung ein Vermögensüberschuss bestehen, geht dieser an
  den Verein der Studierenden der HSR (VSHSR) über.
\lo


\section{Übergangs- und Vollzugsbestimmungen}

\ol
  \li Bei Unklarheiten über die Auslegung dieser Statuten entscheidet der
  Vorstand abschliessend.
\lo


\end{document}
